% =====================================================================
%  minggu-14.tex  --  Slide Pertemuan Minggu ke-14
%  Pengantar Kriptografi (MPFA014) -- Prodi Matematika FMIPA UBB
%
%  Topik silabus : Knapsack, steganografi, dan steganalisis
%  Sumber materi : buku_ajar/chapters/13 (Bab 13) untuk knapsack;
%                  referensi/19-Steganografi-Bagian1-2026 dan
%                  referensi/20-Steganografi-Bagian2-2026 (modul 19 dan 20)
%                  untuk steganografi dan steganalisis
%  Sub-CPMK      : 1.4 dan 2.5
%
%  Kompilasi dari dalam folder ppt/:
%      pdflatex -interaction=nonstopmode minggu-14.tex   (dua kali)
% =====================================================================

\documentclass[aspectratio=169,11pt]{beamer}
% =====================================================================
%  beamer-preamble.tex  --  prasarana dan gaya slide kuliah
%
%  Mata kuliah : Pengantar Kriptografi (MPFA014)
%  Program     : Program Studi Matematika, FMIPA, Universitas Bale Bandung
%
%  Berkas ini TIDAK memuat \documentclass. Setiap berkas minggu memuatnya
%  sendiri lalu menarik berkas ini:
%
%      \documentclass[aspectratio=169,11pt]{beamer}
%      % =====================================================================
%  beamer-preamble.tex  --  prasarana dan gaya slide kuliah
%
%  Mata kuliah : Pengantar Kriptografi (MPFA014)
%  Program     : Program Studi Matematika, FMIPA, Universitas Bale Bandung
%
%  Berkas ini TIDAK memuat \documentclass. Setiap berkas minggu memuatnya
%  sendiri lalu menarik berkas ini:
%
%      \documentclass[aspectratio=169,11pt]{beamer}
%      % =====================================================================
%  beamer-preamble.tex  --  prasarana dan gaya slide kuliah
%
%  Mata kuliah : Pengantar Kriptografi (MPFA014)
%  Program     : Program Studi Matematika, FMIPA, Universitas Bale Bandung
%
%  Berkas ini TIDAK memuat \documentclass. Setiap berkas minggu memuatnya
%  sendiri lalu menarik berkas ini:
%
%      \documentclass[aspectratio=169,11pt]{beamer}
%      \input{beamer-preamble}
%      \title{...} \subtitle{...} \setWeekInfo{Minggu n}{Topik}
%
%  Berkas minggu hanya berisi \setWeekInfo dan lingkungan frame.
%  Kompilasi dijalankan DARI DALAM folder ppt/ (lihat compile-ppt.bat),
%  karena gambar buku dirujuk dengan jalur relatif ../buku_ajar/.
% =====================================================================

% ---------------------------------------------------------------- bahasa
\usepackage[utf8]{inputenc}
\usepackage[T1]{fontenc}
\usepackage[indonesian]{babel}
\usepackage{lmodern}
\usepackage{microtype}

% ------------------------------------------------------------- matematika
\usepackage{amsmath,amssymb,mathtools}

% ----------------------------------------------------------------- tabel
\usepackage{booktabs}
\usepackage{array}
\usepackage{tabularx}

% --------------------------------------------------------------- grafik
\usepackage{tikz}
\usetikzlibrary{arrows.meta,positioning,shapes.geometric,calc,fit,backgrounds}

% ---------------------------------------------------------- kode program
\usepackage{listings}
\lstdefinestyle{slidecode}{
    basicstyle=\ttfamily\scriptsize,
    keywordstyle=\color{ubbblue}\bfseries,
    commentstyle=\color{ubbgreen}\itshape,
    stringstyle=\color{ubbred},
    numbers=left,
    numberstyle=\tiny\color{gray},
    numbersep=6pt,
    frame=single,
    rulecolor=\color{ubbblue!40},
    backgroundcolor=\color{ubbblue!3},
    breaklines=true,
    showstringspaces=false,
    tabsize=4,
    columns=fullflexible,
    keepspaces=true,
}
\lstset{style=slidecode}

% =====================================================================
%  WARNA INSTITUSI
% =====================================================================
\definecolor{ubbblue}{RGB}{12, 62, 120}      % warna utama
\definecolor{ubbteal}{RGB}{0, 118, 110}      % aksen kedua
\definecolor{ubbgreen}{RGB}{22, 122, 74}     % untuk blok contoh
\definecolor{ubbred}{RGB}{168, 42, 42}       % untuk blok peringatan
\definecolor{ubbgray}{RGB}{70, 76, 84}       % teks sekunder

% =====================================================================
%  IDENTITAS MATA KULIAH
% =====================================================================
\newcommand{\CourseName}{Pengantar Kriptografi}
\newcommand{\CourseCode}{MPFA014}
\newcommand{\CourseSemester}{Semester Ganjil}
\newcommand{\CourseProgram}{Program Studi Matematika}
\newcommand{\CourseFaculty}{Fakultas Matematika dan Ilmu Pengetahuan Alam}
\newcommand{\CourseInstitution}{Universitas Bale Bandung}
\newcommand{\CourseLecturer}{Cecep Suwanda, S.Si., M.Kom.}

% Label minggu, disetel tiap berkas lewat \setWeekInfo
\newcommand{\WeekLabel}{}
\newcommand{\WeekTopic}{}
\newcommand{\setWeekInfo}[2]{%
    \renewcommand{\WeekLabel}{#1}%
    \renewcommand{\WeekTopic}{#2}%
}

% =====================================================================
%  TEMA: header, footer, warna blok
% =====================================================================
\usetheme{default}
\setbeamertemplate{navigation symbols}{}
\setbeamertemplate{blocks}[rounded][shadow=false]

\setbeamercolor{palette primary}{bg=ubbblue, fg=white}
\setbeamercolor{palette secondary}{bg=ubbgray, fg=white}
\setbeamercolor{palette tertiary}{bg=ubbteal, fg=white}

\setbeamercolor{structure}{fg=ubbblue}
\setbeamercolor{normal text}{fg=black!92}
\setbeamercolor{alerted text}{fg=ubbred}

\setbeamercolor{block title}{bg=ubbblue, fg=white}
\setbeamercolor{block body}{bg=ubbblue!6, fg=black!92}
\setbeamercolor{block title alerted}{bg=ubbred, fg=white}
\setbeamercolor{block body alerted}{bg=ubbred!6, fg=black!92}
\setbeamercolor{block title example}{bg=ubbgreen, fg=white}
\setbeamercolor{block body example}{bg=ubbgreen!6, fg=black!92}

\setbeamercolor{itemize item}{fg=ubbblue}
\setbeamercolor{itemize subitem}{fg=ubbteal}
\setbeamercolor{enumerate item}{fg=ubbblue}

\setbeamerfont{headline}{size=\scriptsize}
\setbeamerfont{footline}{size=\scriptsize}
\setbeamerfont{block title}{size=\normalsize}

% Badan blok rata kiri (bukan rata kanan-kiri): teks slide lebih mudah dibaca
% dan baris terakhir yang pendek tidak lagi memicu "Underfull \hbox".
\addtobeamertemplate{block body begin}{\raggedright}{}
\addtobeamertemplate{block alerted body begin}{\raggedright}{}
\addtobeamertemplate{block example body begin}{\raggedright}{}

% Header: nama mata kuliah dan kode MK
\setbeamertemplate{headline}{%
    \begin{beamercolorbox}[wd=\paperwidth,ht=2.6ex,dp=1.2ex,leftskip=1em,rightskip=1em]{palette primary}%
        \usebeamerfont{headline}\CourseName\ \textperiodcentered\ \CourseCode%
        \hfill \CourseProgram%
    \end{beamercolorbox}%
}

% Footer: minggu dan nomor halaman
\setbeamertemplate{footline}{%
    \begin{beamercolorbox}[wd=\paperwidth,ht=2.4ex,dp=1.2ex,leftskip=1em,rightskip=1em]{palette secondary}%
        \usebeamerfont{footline}\WeekLabel\ \textperiodcentered\ \WeekTopic%
        \hfill \insertframenumber\,/\,\inserttotalframenumber%
    \end{beamercolorbox}%
}

% Halaman judul: judul, subjudul, lalu blok identitas mata kuliah
\setbeamertemplate{title page}{%
    \begin{centering}
        \vskip1.0em
        {\usebeamerfont{title}\usebeamercolor[fg]{title}\inserttitle\par}
        \vskip0.3em
        {\usebeamerfont{subtitle}\usebeamercolor[fg]{subtitle}\insertsubtitle\par}
        \vskip1.2em
        \textcolor{ubbblue}{\rule{0.72\textwidth}{0.8pt}}
        \vskip1.2em
        {\small
            \textbf{\CourseName}\ (\CourseCode)\par
            \vskip0.4em
            \CourseProgram\par
            \CourseFaculty\par
            \CourseInstitution\par
            \vskip0.6em
            \CourseSemester\ \textperiodcentered\ \WeekLabel\par
            \vskip0.8em
            Dosen Pengampu: \CourseLecturer\par
        }
        \vskip1.0em
        \textcolor{ubbblue}{\rule{0.72\textwidth}{0.4pt}}
    \end{centering}
}

% =====================================================================
%  PERINTAH BANTU
% =====================================================================
\newcommand{\R}{\mathbb{R}}
\newcommand{\N}{\mathbb{N}}
\newcommand{\Z}{\mathbb{Z}}
\newcommand{\code}[1]{\texttt{#1}}

% Rujukan gambar dari buku: \bookfig[lebar]{NN}{namaberkas.png}
% Jalur relatif terhadap folder ppt/ (dikompilasi dari sana).
\newcommand{\bookfig}[3][0.80\linewidth]{%
    \includegraphics[width=#1]{../buku_ajar/chapters/#2/figures/#3}%
}

% Keterangan sumber gambar, sepadan dengan \sumber di buku
\newcommand{\figsumber}[1]{%
    \par\vspace{1pt}{\tiny\itshape Sumber: #1\par}%
}

% Penanda blok materi yang sering dipakai di slide
\newcommand{\slidejudul}[1]{{\large\bfseries\textcolor{ubbblue}{#1}}\par\vspace{2pt}}
\newcommand{\kunci}[1]{\textcolor{ubbred}{\textbf{#1}}}

%      \title{...} \subtitle{...} \setWeekInfo{Minggu n}{Topik}
%
%  Berkas minggu hanya berisi \setWeekInfo dan lingkungan frame.
%  Kompilasi dijalankan DARI DALAM folder ppt/ (lihat compile-ppt.bat),
%  karena gambar buku dirujuk dengan jalur relatif ../buku_ajar/.
% =====================================================================

% ---------------------------------------------------------------- bahasa
\usepackage[utf8]{inputenc}
\usepackage[T1]{fontenc}
\usepackage[indonesian]{babel}
\usepackage{lmodern}
\usepackage{microtype}

% ------------------------------------------------------------- matematika
\usepackage{amsmath,amssymb,mathtools}

% ----------------------------------------------------------------- tabel
\usepackage{booktabs}
\usepackage{array}
\usepackage{tabularx}

% --------------------------------------------------------------- grafik
\usepackage{tikz}
\usetikzlibrary{arrows.meta,positioning,shapes.geometric,calc,fit,backgrounds}

% ---------------------------------------------------------- kode program
\usepackage{listings}
\lstdefinestyle{slidecode}{
    basicstyle=\ttfamily\scriptsize,
    keywordstyle=\color{ubbblue}\bfseries,
    commentstyle=\color{ubbgreen}\itshape,
    stringstyle=\color{ubbred},
    numbers=left,
    numberstyle=\tiny\color{gray},
    numbersep=6pt,
    frame=single,
    rulecolor=\color{ubbblue!40},
    backgroundcolor=\color{ubbblue!3},
    breaklines=true,
    showstringspaces=false,
    tabsize=4,
    columns=fullflexible,
    keepspaces=true,
}
\lstset{style=slidecode}

% =====================================================================
%  WARNA INSTITUSI
% =====================================================================
\definecolor{ubbblue}{RGB}{12, 62, 120}      % warna utama
\definecolor{ubbteal}{RGB}{0, 118, 110}      % aksen kedua
\definecolor{ubbgreen}{RGB}{22, 122, 74}     % untuk blok contoh
\definecolor{ubbred}{RGB}{168, 42, 42}       % untuk blok peringatan
\definecolor{ubbgray}{RGB}{70, 76, 84}       % teks sekunder

% =====================================================================
%  IDENTITAS MATA KULIAH
% =====================================================================
\newcommand{\CourseName}{Pengantar Kriptografi}
\newcommand{\CourseCode}{MPFA014}
\newcommand{\CourseSemester}{Semester Ganjil}
\newcommand{\CourseProgram}{Program Studi Matematika}
\newcommand{\CourseFaculty}{Fakultas Matematika dan Ilmu Pengetahuan Alam}
\newcommand{\CourseInstitution}{Universitas Bale Bandung}
\newcommand{\CourseLecturer}{Cecep Suwanda, S.Si., M.Kom.}

% Label minggu, disetel tiap berkas lewat \setWeekInfo
\newcommand{\WeekLabel}{}
\newcommand{\WeekTopic}{}
\newcommand{\setWeekInfo}[2]{%
    \renewcommand{\WeekLabel}{#1}%
    \renewcommand{\WeekTopic}{#2}%
}

% =====================================================================
%  TEMA: header, footer, warna blok
% =====================================================================
\usetheme{default}
\setbeamertemplate{navigation symbols}{}
\setbeamertemplate{blocks}[rounded][shadow=false]

\setbeamercolor{palette primary}{bg=ubbblue, fg=white}
\setbeamercolor{palette secondary}{bg=ubbgray, fg=white}
\setbeamercolor{palette tertiary}{bg=ubbteal, fg=white}

\setbeamercolor{structure}{fg=ubbblue}
\setbeamercolor{normal text}{fg=black!92}
\setbeamercolor{alerted text}{fg=ubbred}

\setbeamercolor{block title}{bg=ubbblue, fg=white}
\setbeamercolor{block body}{bg=ubbblue!6, fg=black!92}
\setbeamercolor{block title alerted}{bg=ubbred, fg=white}
\setbeamercolor{block body alerted}{bg=ubbred!6, fg=black!92}
\setbeamercolor{block title example}{bg=ubbgreen, fg=white}
\setbeamercolor{block body example}{bg=ubbgreen!6, fg=black!92}

\setbeamercolor{itemize item}{fg=ubbblue}
\setbeamercolor{itemize subitem}{fg=ubbteal}
\setbeamercolor{enumerate item}{fg=ubbblue}

\setbeamerfont{headline}{size=\scriptsize}
\setbeamerfont{footline}{size=\scriptsize}
\setbeamerfont{block title}{size=\normalsize}

% Badan blok rata kiri (bukan rata kanan-kiri): teks slide lebih mudah dibaca
% dan baris terakhir yang pendek tidak lagi memicu "Underfull \hbox".
\addtobeamertemplate{block body begin}{\raggedright}{}
\addtobeamertemplate{block alerted body begin}{\raggedright}{}
\addtobeamertemplate{block example body begin}{\raggedright}{}

% Header: nama mata kuliah dan kode MK
\setbeamertemplate{headline}{%
    \begin{beamercolorbox}[wd=\paperwidth,ht=2.6ex,dp=1.2ex,leftskip=1em,rightskip=1em]{palette primary}%
        \usebeamerfont{headline}\CourseName\ \textperiodcentered\ \CourseCode%
        \hfill \CourseProgram%
    \end{beamercolorbox}%
}

% Footer: minggu dan nomor halaman
\setbeamertemplate{footline}{%
    \begin{beamercolorbox}[wd=\paperwidth,ht=2.4ex,dp=1.2ex,leftskip=1em,rightskip=1em]{palette secondary}%
        \usebeamerfont{footline}\WeekLabel\ \textperiodcentered\ \WeekTopic%
        \hfill \insertframenumber\,/\,\inserttotalframenumber%
    \end{beamercolorbox}%
}

% Halaman judul: judul, subjudul, lalu blok identitas mata kuliah
\setbeamertemplate{title page}{%
    \begin{centering}
        \vskip1.0em
        {\usebeamerfont{title}\usebeamercolor[fg]{title}\inserttitle\par}
        \vskip0.3em
        {\usebeamerfont{subtitle}\usebeamercolor[fg]{subtitle}\insertsubtitle\par}
        \vskip1.2em
        \textcolor{ubbblue}{\rule{0.72\textwidth}{0.8pt}}
        \vskip1.2em
        {\small
            \textbf{\CourseName}\ (\CourseCode)\par
            \vskip0.4em
            \CourseProgram\par
            \CourseFaculty\par
            \CourseInstitution\par
            \vskip0.6em
            \CourseSemester\ \textperiodcentered\ \WeekLabel\par
            \vskip0.8em
            Dosen Pengampu: \CourseLecturer\par
        }
        \vskip1.0em
        \textcolor{ubbblue}{\rule{0.72\textwidth}{0.4pt}}
    \end{centering}
}

% =====================================================================
%  PERINTAH BANTU
% =====================================================================
\newcommand{\R}{\mathbb{R}}
\newcommand{\N}{\mathbb{N}}
\newcommand{\Z}{\mathbb{Z}}
\newcommand{\code}[1]{\texttt{#1}}

% Rujukan gambar dari buku: \bookfig[lebar]{NN}{namaberkas.png}
% Jalur relatif terhadap folder ppt/ (dikompilasi dari sana).
\newcommand{\bookfig}[3][0.80\linewidth]{%
    \includegraphics[width=#1]{../buku_ajar/chapters/#2/figures/#3}%
}

% Keterangan sumber gambar, sepadan dengan \sumber di buku
\newcommand{\figsumber}[1]{%
    \par\vspace{1pt}{\tiny\itshape Sumber: #1\par}%
}

% Penanda blok materi yang sering dipakai di slide
\newcommand{\slidejudul}[1]{{\large\bfseries\textcolor{ubbblue}{#1}}\par\vspace{2pt}}
\newcommand{\kunci}[1]{\textcolor{ubbred}{\textbf{#1}}}

%      \title{...} \subtitle{...} \setWeekInfo{Minggu n}{Topik}
%
%  Berkas minggu hanya berisi \setWeekInfo dan lingkungan frame.
%  Kompilasi dijalankan DARI DALAM folder ppt/ (lihat compile-ppt.bat),
%  karena gambar buku dirujuk dengan jalur relatif ../buku_ajar/.
% =====================================================================

% ---------------------------------------------------------------- bahasa
\usepackage[utf8]{inputenc}
\usepackage[T1]{fontenc}
\usepackage[indonesian]{babel}
\usepackage{lmodern}
\usepackage{microtype}

% ------------------------------------------------------------- matematika
\usepackage{amsmath,amssymb,mathtools}

% ----------------------------------------------------------------- tabel
\usepackage{booktabs}
\usepackage{array}
\usepackage{tabularx}

% --------------------------------------------------------------- grafik
\usepackage{tikz}
\usetikzlibrary{arrows.meta,positioning,shapes.geometric,calc,fit,backgrounds}

% ---------------------------------------------------------- kode program
\usepackage{listings}
\lstdefinestyle{slidecode}{
    basicstyle=\ttfamily\scriptsize,
    keywordstyle=\color{ubbblue}\bfseries,
    commentstyle=\color{ubbgreen}\itshape,
    stringstyle=\color{ubbred},
    numbers=left,
    numberstyle=\tiny\color{gray},
    numbersep=6pt,
    frame=single,
    rulecolor=\color{ubbblue!40},
    backgroundcolor=\color{ubbblue!3},
    breaklines=true,
    showstringspaces=false,
    tabsize=4,
    columns=fullflexible,
    keepspaces=true,
}
\lstset{style=slidecode}

% =====================================================================
%  WARNA INSTITUSI
% =====================================================================
\definecolor{ubbblue}{RGB}{12, 62, 120}      % warna utama
\definecolor{ubbteal}{RGB}{0, 118, 110}      % aksen kedua
\definecolor{ubbgreen}{RGB}{22, 122, 74}     % untuk blok contoh
\definecolor{ubbred}{RGB}{168, 42, 42}       % untuk blok peringatan
\definecolor{ubbgray}{RGB}{70, 76, 84}       % teks sekunder

% =====================================================================
%  IDENTITAS MATA KULIAH
% =====================================================================
\newcommand{\CourseName}{Pengantar Kriptografi}
\newcommand{\CourseCode}{MPFA014}
\newcommand{\CourseSemester}{Semester Ganjil}
\newcommand{\CourseProgram}{Program Studi Matematika}
\newcommand{\CourseFaculty}{Fakultas Matematika dan Ilmu Pengetahuan Alam}
\newcommand{\CourseInstitution}{Universitas Bale Bandung}
\newcommand{\CourseLecturer}{Cecep Suwanda, S.Si., M.Kom.}

% Label minggu, disetel tiap berkas lewat \setWeekInfo
\newcommand{\WeekLabel}{}
\newcommand{\WeekTopic}{}
\newcommand{\setWeekInfo}[2]{%
    \renewcommand{\WeekLabel}{#1}%
    \renewcommand{\WeekTopic}{#2}%
}

% =====================================================================
%  TEMA: header, footer, warna blok
% =====================================================================
\usetheme{default}
\setbeamertemplate{navigation symbols}{}
\setbeamertemplate{blocks}[rounded][shadow=false]

\setbeamercolor{palette primary}{bg=ubbblue, fg=white}
\setbeamercolor{palette secondary}{bg=ubbgray, fg=white}
\setbeamercolor{palette tertiary}{bg=ubbteal, fg=white}

\setbeamercolor{structure}{fg=ubbblue}
\setbeamercolor{normal text}{fg=black!92}
\setbeamercolor{alerted text}{fg=ubbred}

\setbeamercolor{block title}{bg=ubbblue, fg=white}
\setbeamercolor{block body}{bg=ubbblue!6, fg=black!92}
\setbeamercolor{block title alerted}{bg=ubbred, fg=white}
\setbeamercolor{block body alerted}{bg=ubbred!6, fg=black!92}
\setbeamercolor{block title example}{bg=ubbgreen, fg=white}
\setbeamercolor{block body example}{bg=ubbgreen!6, fg=black!92}

\setbeamercolor{itemize item}{fg=ubbblue}
\setbeamercolor{itemize subitem}{fg=ubbteal}
\setbeamercolor{enumerate item}{fg=ubbblue}

\setbeamerfont{headline}{size=\scriptsize}
\setbeamerfont{footline}{size=\scriptsize}
\setbeamerfont{block title}{size=\normalsize}

% Badan blok rata kiri (bukan rata kanan-kiri): teks slide lebih mudah dibaca
% dan baris terakhir yang pendek tidak lagi memicu "Underfull \hbox".
\addtobeamertemplate{block body begin}{\raggedright}{}
\addtobeamertemplate{block alerted body begin}{\raggedright}{}
\addtobeamertemplate{block example body begin}{\raggedright}{}

% Header: nama mata kuliah dan kode MK
\setbeamertemplate{headline}{%
    \begin{beamercolorbox}[wd=\paperwidth,ht=2.6ex,dp=1.2ex,leftskip=1em,rightskip=1em]{palette primary}%
        \usebeamerfont{headline}\CourseName\ \textperiodcentered\ \CourseCode%
        \hfill \CourseProgram%
    \end{beamercolorbox}%
}

% Footer: minggu dan nomor halaman
\setbeamertemplate{footline}{%
    \begin{beamercolorbox}[wd=\paperwidth,ht=2.4ex,dp=1.2ex,leftskip=1em,rightskip=1em]{palette secondary}%
        \usebeamerfont{footline}\WeekLabel\ \textperiodcentered\ \WeekTopic%
        \hfill \insertframenumber\,/\,\inserttotalframenumber%
    \end{beamercolorbox}%
}

% Halaman judul: judul, subjudul, lalu blok identitas mata kuliah
\setbeamertemplate{title page}{%
    \begin{centering}
        \vskip1.0em
        {\usebeamerfont{title}\usebeamercolor[fg]{title}\inserttitle\par}
        \vskip0.3em
        {\usebeamerfont{subtitle}\usebeamercolor[fg]{subtitle}\insertsubtitle\par}
        \vskip1.2em
        \textcolor{ubbblue}{\rule{0.72\textwidth}{0.8pt}}
        \vskip1.2em
        {\small
            \textbf{\CourseName}\ (\CourseCode)\par
            \vskip0.4em
            \CourseProgram\par
            \CourseFaculty\par
            \CourseInstitution\par
            \vskip0.6em
            \CourseSemester\ \textperiodcentered\ \WeekLabel\par
            \vskip0.8em
            Dosen Pengampu: \CourseLecturer\par
        }
        \vskip1.0em
        \textcolor{ubbblue}{\rule{0.72\textwidth}{0.4pt}}
    \end{centering}
}

% =====================================================================
%  PERINTAH BANTU
% =====================================================================
\newcommand{\R}{\mathbb{R}}
\newcommand{\N}{\mathbb{N}}
\newcommand{\Z}{\mathbb{Z}}
\newcommand{\code}[1]{\texttt{#1}}

% Rujukan gambar dari buku: \bookfig[lebar]{NN}{namaberkas.png}
% Jalur relatif terhadap folder ppt/ (dikompilasi dari sana).
\newcommand{\bookfig}[3][0.80\linewidth]{%
    \includegraphics[width=#1]{../buku_ajar/chapters/#2/figures/#3}%
}

% Keterangan sumber gambar, sepadan dengan \sumber di buku
\newcommand{\figsumber}[1]{%
    \par\vspace{1pt}{\tiny\itshape Sumber: #1\par}%
}

% Penanda blok materi yang sering dipakai di slide
\newcommand{\slidejudul}[1]{{\large\bfseries\textcolor{ubbblue}{#1}}\par\vspace{2pt}}
\newcommand{\kunci}[1]{\textcolor{ubbred}{\textbf{#1}}}


\title{Knapsack dan Steganografi: Dua Jalan Menyembunyikan}
\subtitle{Sistem Merkle--Hellman, LSB pada citra dan audio, serta steganalisis}
\setWeekInfo{Minggu 14}{Knapsack, steganografi, dan steganalisis}

\begin{document}

% ----------------------------------------------------------- halaman judul
\begin{frame}[plain]
    \titlepage
\end{frame}

% ------------------------------------------------------------------ peta
\begin{frame}{Peta Materi Pertemuan Ini}
    \begin{columns}[T,onlytextwidth]
        \begin{column}{0.54\textwidth}
            \slidejudul{Lima bagian}
            \footnotesize
            \begin{enumerate}
                \item \textbf{Persoalan knapsack} --- memilih objek dengan
                      bobot total tertentu.
                \item \textbf{Barisan \textit{superincreasing} dan
                      Merkle--Hellman} --- kunci privat yang mudah, kunci
                      publik yang tampak sukar.
                \item \textbf{Contoh terhitung} --- pembangkitan kunci,
                      enkripsi, dan dekripsi dengan angka dari bab.
                \item \textbf{Steganografi} --- gagasan penyembunyian
                      informasi dan LSB pada citra serta audio.
                \item \textbf{Steganalisis} --- tujuan, jenis, dan metode
                      pendeteksian pesan tersembunyi.
            \end{enumerate}
        \end{column}
        \begin{column}{0.42\textwidth}
            \begin{block}{Referensi}
                \footnotesize
                Bab 13 buku ajar; Munir, \textit{IF4020 Kriptografi}, ITB,
                2026 (modul 16, 19, dan 20)
            \end{block}
            \begin{exampleblock}{Benang merah}
                \footnotesize
                Knapsack menyembunyikan pesan di balik \kunci{kesulitan
                hitung}; steganografi menyembunyikan pesan di balik
                \kunci{media yang tampak biasa}.
            \end{exampleblock}
        \end{column}
    \end{columns}
\end{frame}

% ---------------------------------------------------- capaian pembelajaran
\begin{frame}{Capaian Pembelajaran}
    \begin{block}{Sub-CPMK yang dibebankan}
        \footnotesize
        \begin{itemize}
            \item \textbf{Sub-CPMK 1.4} --- Mengenal algoritma kriptografi
                  kunci publik, yaitu RSA, Diffie--Hellman, ElGamal, knapsack,
                  dan \textit{Elliptic Curve Cryptography} (ECC), beserta
                  persoalan matematis yang mendasarinya.
            \item \textbf{Sub-CPMK 2.5} --- Menjelaskan gagasan steganografi
                  dan steganalisis sebagai pelengkap kriptografi dalam
                  penyembunyian informasi beserta urgensinya.
        \end{itemize}
    \end{block}
    \vspace{0.1cm}
    \begin{exampleblock}{Setelah pertemuan ini, mahasiswa mampu}
        \footnotesize
        \begin{itemize}
            \item menjelaskan persoalan knapsack dan barisan
                  \textit{superincreasing};
            \item menjelaskan prosedur Merkle--Hellman dari pembangkitan kunci
                  sampai dekripsi;
            \item menghitung satu contoh enkripsi dan dekripsi knapsack;
            \item menjelaskan mengapa knapsack tidak lagi dianggap aman;
            \item menjelaskan LSB pada citra dan audio; serta
            \item membedakan jenis dan metode steganalisis.
        \end{itemize}
    \end{exampleblock}
\end{frame}

% =====================================================================
\section{Knapsack: dari Tas ke Kunci Publik}
% =====================================================================

\begin{frame}{Persoalan Knapsack}
    Diberikan $n$ objek dengan bobot $w_1, w_2, \dots, w_n$ dan sebuah tas
    berkapasitas total $M$. Kita harus memilih objek yang dimasukkan:
    \[
        M = b_1w_1 + b_2w_2 + \cdots + b_nw_n,
        \qquad b_i \in \{0,1\}.
    \]
    \begin{itemize}
        \item $b_i = 1$ berarti objek ke-$i$ dimasukkan; $b_i = 0$ berarti
              tidak.
        \item Persoalan ini termasuk kelompok \kunci{NP-complete}.
        \item Tidak ada algoritma waktu polinomial yang diketahui untuk
              memecahkannya ketika $n$ membesar.
    \end{itemize}
    \vspace{0.15cm}
    \begin{exampleblock}{Asimetri yang dimanfaatkan}
        \footnotesize
        Menghitung $M$ dari bit $b_i$ itu mudah. Menentukan kembali bit $b_i$
        dari $M$ itu sukar.
    \end{exampleblock}
\end{frame}

\begin{frame}{Ilustrasi Persoalan Knapsack}
    \begin{columns}[T,onlytextwidth]
        \begin{column}{0.50\textwidth}
            \centering
            \bookfig[0.95\linewidth]{13}{ilustrasi-masalah-knapsack.png}
            \figsumber{Munir, \textit{IF4020 Kriptografi}, ITB.}
        \end{column}
        \begin{column}{0.46\textwidth}
            \footnotesize
            \begin{itemize}
                \item Setiap barang memiliki bobot yang berbeda.
                \item Kita harus memilih himpunan bagian yang berat totalnya
                      tepat $15$ kg.
                \item Dalam kriptografi, pilihan itu menyatakan bit $0$ atau
                      $1$.
                \item Menghitung total itu mudah; menebak pilihan yang
                      menghasilkan total itu yang sukar.
            \end{itemize}
            \begin{exampleblock}{Inti persoalannya}
                \footnotesize
                Banyak himpunan bagian yang mungkin, tetapi hanya sedikit yang
                memberikan total tepat.
            \end{exampleblock}
        \end{column}
    \end{columns}
\end{frame}

\begin{frame}{Dari Persoalan ke Kriptografi}
    \begin{columns}[T,onlytextwidth]
        \begin{column}{0.50\textwidth}
            \footnotesize
            \begin{itemize}
                \item Bit-bit plainteks berperan sebagai $b_i$.
                \item Bobot $w_i$ berperan sebagai kunci rahasia.
                \item Pengirim menghasilkan satu bilangan $M$ untuk setiap
                      blok bit.
                \item Penyadap harus menyelesaikan persoalan subset-sum untuk
                      membacanya.
            \end{itemize}
        \end{column}
        \begin{column}{0.46\textwidth}
            \begin{alertblock}{Masalah pertama}
                \footnotesize
                Penerima yang sah juga tidak dapat menyelesaikan persoalan itu.
                Knapsack biasa hanya dapat dipakai untuk enkripsi, bukan untuk
                dekripsi.
            \end{alertblock}
            \begin{block}{Jalan keluar}
                \footnotesize
                Pilih barisan bobot yang \kunci{mudah dipecahkan oleh pemilik
                kunci}, tetapi tampak sukar bagi orang lain.
            \end{block}
        \end{column}
    \end{columns}
\end{frame}

\begin{frame}{Barisan \textit{Superincreasing}}
    Barisan $\{w_1, w_2, \dots, w_n\}$ disebut \textit{superincreasing} jika
    setiap elemen lebih besar daripada jumlah semua elemen sebelumnya:
    \[
        w_i > \sum_{j=1}^{i-1} w_j
        \quad \text{untuk } i = 2, \dots, n.
    \]
    \begin{columns}[T,onlytextwidth]
        \begin{column}{0.50\textwidth}
            \begin{exampleblock}{Contoh yang memenuhi}
                \footnotesize
                $\{1,3,6,13,27,52\}$:
                $3>1$, $6>4$, $13>10$, $27>23$, dan $52>50$.
            \end{exampleblock}
        \end{column}
        \begin{column}{0.46\textwidth}
            \begin{alertblock}{Contoh yang gagal}
                \footnotesize
                $\{1,3,4,9,15,25\}$:
                $4$ tidak lebih besar daripada $1+3=4$, dan $25$ tidak lebih
                besar daripada $1+3+4+9+15=32$.
            \end{alertblock}
        \end{column}
    \end{columns}
    \vspace{0.05cm}
    \footnotesize
    Akibatnya, elemen terbesar tidak mungkin digantikan oleh kombinasi
    elemen-elemen yang lebih kecil.
\end{frame}

\begin{frame}{Prosedur \textit{Greedy} pada Barisan \textit{Superincreasing}}
    \begin{enumerate}
        \item Mulai dengan bobot total $M$.
        \item Bandingkan bobot total dengan bobot terbesar yang tersisa.
        \item Jika bobot itu $\le$ bobot total, masukkan objeknya dan kurangi
              bobot total. Jika tidak, lewati objek itu.
        \item Ulangi untuk bobot berikutnya sampai semua bobot diperiksa.
        \item Jika bobot total menjadi nol, solusinya ditemukan.
    \end{enumerate}
    \vspace{0.15cm}
    \begin{columns}[T,onlytextwidth]
        \begin{column}{0.48\textwidth}
            \begin{block}{Mengapa selalu berhasil}
                \footnotesize
                Keputusan untuk elemen terbesar bersifat pasti: elemen itu
                tidak dapat dibentuk oleh semua elemen yang lebih kecil.
            \end{block}
        \end{column}
        \begin{column}{0.48\textwidth}
            \begin{alertblock}{Harga kemudahan}
                \footnotesize
                Dekripsi berlangsung dalam waktu \kunci{linear}, $O(n)$.
                Siapa pun yang mengetahui barisannya dapat membuka pesan.
            \end{alertblock}
        \end{column}
    \end{columns}
\end{frame}

\begin{frame}{Gagasan Merkle--Hellman}
    \begin{columns}[T,onlytextwidth]
        \begin{column}{0.54\textwidth}
            \footnotesize
            \begin{itemize}
                \item Merkle dan Hellman (1978) menyembunyikan barisan
                      \textit{superincreasing} di balik perkalian modulo.
                \item Kunci privat tetap barisan \textit{superincreasing}
                      $W_{priv}$.
                \item Kunci publik $W_{pub}$ dibentuk dengan mengalikan setiap
                      elemen lalu mengambil modulo:
                      \[
                          w'_i = (w_i \cdot n) \bmod m.
                      \]
                \item Hasilnya bukan lagi barisan \textit{superincreasing}.
            \end{itemize}
        \end{column}
        \begin{column}{0.42\textwidth}
            \begin{exampleblock}{Dua kunci, dua peran}
                \footnotesize
                Pengirim memakai $W_{pub}$ untuk enkripsi. Penerima memakai
                $W_{priv}$ dan $(n,m)$ untuk dekripsi.
            \end{exampleblock}
            \begin{alertblock}{Catatan notasi}
                \footnotesize
                $n$ di sini adalah \kunci{pengali}, bukan jumlah elemen.
                Jumlah elemen ditulis melalui indeks $i$.
            \end{alertblock}
        \end{column}
    \end{columns}
\end{frame}

\begin{frame}{Syarat Pembangkitan Kunci}
    \begin{block}{Dua syarat yang tidak boleh dilanggar}
        \begin{enumerate}
            \item Modulus harus lebih besar daripada jumlah seluruh elemen
                  privat:
                  \[
                      m > \sum_{i=1}^{n} w_i.
                  \]
            \item Pengali $n$ harus relatif prima terhadap modulus:
                  \[
                      \mathrm{PBB}(n,m) = 1.
                  \]
        \end{enumerate}
    \end{block}
    \vspace{0.15cm}
    \begin{columns}[T,onlytextwidth]
        \begin{column}{0.48\textwidth}
            \begin{exampleblock}{Akibat syarat pertama}
                \footnotesize
                Penjumlahan seluruh elemen privat tidak pernah melewati
                modulus, sehingga hasil dekripsi berupa kesamaan biasa.
            \end{exampleblock}
        \end{column}
        \begin{column}{0.48\textwidth}
            \begin{alertblock}{Akibat syarat kedua}
                \footnotesize
                $n$ mempunyai invers modulo $m$. Tanpa invers itu, penerima
                yang sah tidak dapat mendekripsi.
            \end{alertblock}
        \end{column}
    \end{columns}
\end{frame}

\begin{frame}{Enkripsi}
    Plainteks dipecah menjadi blok yang panjangnya sama dengan banyaknya
    elemen kunci publik. Untuk setiap blok:
    \[
        C = \sum_{i=1}^{n} b_i w'_i.
    \]
    \begin{itemize}
        \item Bit $1$ menyumbangkan elemen publik yang bersesuaian.
        \item Bit $0$ tidak menyumbangkan apa pun.
        \item Hasilnya satu bilangan $C$ untuk setiap blok.
        \item Penjumlahan \kunci{tidak diambil modulo} $m$.
    \end{itemize}
    \vspace{0.1cm}
    \begin{exampleblock}{Akibatnya}
        \footnotesize
        $C$ dapat lebih besar daripada $m$. Hal itu tidak menjadi masalah
        karena dekripsi hanya memakai $C \bmod m$.
    \end{exampleblock}
\end{frame}

\begin{frame}{Dekripsi}
    Penerima memegang $W_{priv}$, pengali $n$, dan modulus $m$.
    \begin{enumerate}
        \item Hitung invers $n^{-1}$ modulo $m$, yaitu bilangan yang memenuhi
              \[
                  n \cdot n^{-1} \equiv 1 \pmod{m}.
              \]
        \item Kalikan kriptogram dengan invers itu:
              \[
                  C' = (C \cdot n^{-1}) \bmod m.
              \]
        \item Nyatakan $C'$ sebagai jumlah elemen $W_{priv}$ dengan prosedur
              \textit{greedy}.
        \item Bit-bit yang dihasilkan adalah plainteks semula.
    \end{enumerate}
    \vspace{0.1cm}
    \begin{alertblock}{Inti gagasannya}
        \footnotesize
        Penerima tidak perlu memecahkan knapsack publik. Ia mengembalikan
        kriptogram ke ranah privat, lalu memecahkannya secara linear.
    \end{alertblock}
\end{frame}

\begin{frame}{Mengapa Dekripsi Berhasil}
    Substitusikan rumus enkripsi dan pembangkitan kunci:
    \begin{align*}
        C' &\equiv n^{-1}\sum_{i=1}^{n} b_i w'_i \pmod{m} \\
           &\equiv n^{-1}\sum_{i=1}^{n} b_i (w_i n) \pmod{m} \\
           &\equiv \sum_{i=1}^{n} b_i w_i \pmod{m}.
    \end{align*}
    \begin{itemize}
        \item Faktor $n^{-1}$ dan $n$ saling meniadakan.
        \item Karena $0 \le \sum b_iw_i \le \sum w_i < m$, kongruensi itu
              menjadi \kunci{kesamaan biasa}.
        \item Jadi $C' = \sum b_iw_i$, dan \textit{greedy} memulihkan bit
              dengan tepat.
    \end{itemize}
    \vspace{0.05cm}
    \footnotesize
    Keamanan sistem bergantung pada sulitnya menemukan $n^{-1}$ tanpa
    mengetahui $n$.
\end{frame}

\begin{frame}{Kriptanalisis: Serangan Shamir}
    \begin{itemize}
        \item Pada 1982, Adi Shamir merumuskan transformasi waktu polinomial
              untuk memecahkan Merkle--Hellman dasar.
        \item Shamir tidak memecahkan persoalan knapsack secara umum.
        \item Ia memanfaatkan \kunci{struktur khusus} kunci publik: kunci itu
              merupakan kelipatan modulo dari sebuah barisan
              \textit{superincreasing}.
        \item Struktur tersebut cukup untuk merekonstruksi kunci privat tanpa
              menelusuri seluruh kemungkinan.
    \end{itemize}
    \vspace{0.2cm}
    \begin{alertblock}{Kesimpulan}
        \footnotesize
        Knapsack Merkle--Hellman tidak lagi dianggap aman. Persoalan knapsack
        umum tetap sukar, tetapi \emph{konstruksi khusus} inilah yang runtuh.
    \end{alertblock}
\end{frame}

\begin{frame}{Ukuran Kunci dan Kelemahan Praktis}
    \begin{columns}[T,onlytextwidth]
        \begin{column}{0.50\textwidth}
            \footnotesize
            \begin{itemize}
                \item Kriptogram lebih panjang daripada plainteks: satu blok
                      $n$ bit menjadi satu bilangan besar.
                \item Sumber menganjurkan sedikitnya $250$ elemen.
                \item Setiap elemen sebaiknya berukuran $200$ sampai $400$
                      bit.
            \end{itemize}
        \end{column}
        \begin{column}{0.46\textwidth}
            \begin{alertblock}{Anjuran yang tidak konsisten}
                \footnotesize
                Sumber juga menyebut modulus $100$--$200$ bit. Untuk $250$
                elemen berukuran sedikitnya $200$ bit, jumlahnya sudah melebihi
                $2^{448}$.
            \end{alertblock}
        \end{column}
    \end{columns}
    \vspace{0.15cm}
    \begin{block}{Koreksi yang harus dipahami}
        \footnotesize
        Karena modulus harus lebih besar daripada jumlah seluruh elemen,
        panjangnya paling sedikit \kunci{$449$ bit}. Modulus $100$--$200$ bit
        tidak mungkin memenuhi syarat $m > \sum w_i$.
    \end{block}
\end{frame}

% =====================================================================
\section{Contoh Terhitung Merkle--Hellman}
% =====================================================================

\begin{frame}{Langkah 1: Menyiapkan Kunci Privat}
    Dipilih barisan \textit{superincreasing}
    \[
        W_{priv} = \{2, 3, 6, 13, 27, 52\},
        \qquad \sum_{i=1}^{6} w_i = 103.
    \]
    Pemeriksaan:
    \[
        3>2,\quad 6>2+3=5,\quad 13>2+3+6=11,
    \]
    \[
        27>2+3+6+13=24,\quad 52>2+3+6+13+27=51.
    \]
    \begin{itemize}
        \item Modulus dipilih $m = 105$, lebih besar daripada $103$.
        \item Pengali dipilih $n = 31$.
        \item Karena $\mathrm{PBB}(31,105)=1$, invers $31$ modulo $105$ ada.
    \end{itemize}
\end{frame}

\begin{frame}{Langkah 2: Membangkitkan Kunci Publik}
    Setiap elemen privat dikalikan $31$, lalu diambil modulo $105$:
    \vspace{0.1cm}
    \centering
    \footnotesize
    \begin{tabular}{@{}c r r r@{}}
        \toprule
        $i$ & $w_i$ & $w_i \cdot 31$ & $(w_i \cdot 31) \bmod 105$ \\
        \midrule
        $1$ & $2$  & $62$   & $62$  \\
        $2$ & $3$  & $93$   & $93$  \\
        $3$ & $6$  & $186$  & $81$  \\
        $4$ & $13$ & $403$  & $88$  \\
        $5$ & $27$ & $837$  & $102$ \\
        $6$ & $52$ & $1612$ & $37$  \\
        \bottomrule
    \end{tabular}
    \vspace{0.15cm}
    \raggedright
    \begin{exampleblock}{Hasil}
        \footnotesize
        $W_{pub} = \{62, 93, 81, 88, 102, 37\}$. Barisan ini tidak lagi
        \textit{superincreasing}: $81 < 62+93$, dan $37$ yang tadinya elemen
        terbesar justru menjadi terkecil.
    \end{exampleblock}
\end{frame}

\begin{frame}{Langkah 3: Invers Pengali}
    Penerima membutuhkan $n^{-1} = 31^{-1} \bmod 105$.
    \begin{align*}
        31 \cdot 61 &= 1891, \\
        1891 &= 18 \cdot 105 + 1.
    \end{align*}
    Karena sisanya $1$:
    \[
        n^{-1} = 31^{-1} \bmod 105 = \kunci{61}.
    \]
    \begin{itemize}
        \item Nilai $61$ dipegang penerima dan menjadi rahasia.
        \item Penyerang yang hanya melihat $W_{pub}$ tidak memperoleh $61$.
        \item Pasangan $(n,m)=(31,105)$ dan barisan $W_{priv}$ membentuk
              kunci privat.
    \end{itemize}
    \vspace{0.1cm}
    \begin{exampleblock}{Pesan yang akan dikirim}
        \footnotesize
        Plainteks \code{011000110101101110}, dipecah menjadi tiga blok enam
        bit: \code{011000}, \code{110101}, dan \code{101110}.
    \end{exampleblock}
\end{frame}

\begin{frame}{Langkah 4: Enkripsi}
    Pengirim memakai $W_{pub} = \{62,93,81,88,102,37\}$:
    \vspace{0.1cm}
    \centering
    \footnotesize
    \begin{tabular}{@{}l l r@{}}
        \toprule
        Blok & Bobot publik yang dipakai & Kriptogram $C$ \\
        \midrule
        \code{011000} & $93 + 81$             & $174$ \\
        \code{110101} & $62 + 93 + 88 + 37$   & $280$ \\
        \code{101110} & $62 + 81 + 88 + 102$  & $333$ \\
        \bottomrule
    \end{tabular}
    \vspace{0.15cm}
    \raggedright
    \begin{itemize}
        \item Kriptogram yang dikirim adalah $(174, 280, 333)$.
        \item Nilai $280$ dan $333$ melampaui $m=105$.
        \item Hal itu sah karena dekripsi hanya memakai $C \bmod 105$.
    \end{itemize}
\end{frame}

\begin{frame}{Langkah 5: Dekripsi}
    Setiap kriptogram dikalikan $61$ modulo $105$:
    \vspace{0.1cm}
    \centering
    \footnotesize
    \begin{tabular}{@{}r r r l l@{}}
        \toprule
        $C$ & $C \cdot 61$ & $C' = (C\cdot61)\bmod105$
            & Hasil \textit{greedy} & Blok \\
        \midrule
        $174$ & $10614$ & $9$  & $6+3$             & \code{011000} \\
        $280$ & $17080$ & $70$ & $52+13+3+2$       & \code{110101} \\
        $333$ & $20313$ & $48$ & $27+13+6+2$       & \code{101110} \\
        \bottomrule
    \end{tabular}
    \vspace{0.15cm}
    \raggedright
    \begin{block}{Operasi modulo yang dipakai}
        \footnotesize
        $10614 = 101\cdot105 + 9$;\quad
        $17080 = 162\cdot105 + 70$;\quad
        $20313 = 193\cdot105 + 48$.
    \end{block}
\end{frame}

\begin{frame}{Rincian \textit{Greedy} untuk $C' = 70$}
    \centering
    \footnotesize
    \begin{tabular}{@{}c c c l@{}}
        \toprule
        Bobot & Sisa $M$ & Perbandingan & Bit \\
        \midrule
        $52$ & $70$ & $52 \le 70$ & $b_6 = 1$, sisa $18$ \\
        $27$ & $18$ & $27 > 18$  & $b_5 = 0$ \\
        $13$ & $18$ & $13 \le 18$ & $b_4 = 1$, sisa $5$ \\
        $6$  & $5$  & $6 > 5$    & $b_3 = 0$ \\
        $3$  & $5$  & $3 \le 5$  & $b_2 = 1$, sisa $2$ \\
        $2$  & $2$  & $2 \le 2$  & $b_1 = 1$, sisa $0$ \\
        \bottomrule
    \end{tabular}
    \vspace{0.15cm}
    \raggedright
    \begin{exampleblock}{Membaca hasilnya}
        \footnotesize
        Bit dibaca dari belakang:
        $(b_1,b_2,b_3,b_4,b_5,b_6)=(1,1,0,1,0,1)=\code{110101}$.
        Pemeriksaan: $2+3+13+52=70$.
    \end{exampleblock}
\end{frame}

\begin{frame}{Langkah 6: Memeriksa Hasil}
    Ketiga blok pulih tepat seperti semula:\par
    \vspace{0.1cm}
    \centering
    \footnotesize
    \begin{tabular}{@{}r l l@{}}
        \toprule
        $C'$ & \textit{Greedy} & Blok \\
        \midrule
        $9$  & $6+3$           & \code{011000} \\
        $70$ & $52+13+3+2$     & \code{110101} \\
        $48$ & $27+13+6+2$     & \code{101110} \\
        \bottomrule
    \end{tabular}
    \vspace{0.2cm}
    \raggedright
    \begin{block}{Plainteks yang dipulihkan}
        \footnotesize
        \[
            \code{011000}\;\code{110101}\;\code{101110}
            = \code{011000110101101110}.
        \]
    \end{block}
    \begin{itemize}
        \item Enkripsi hanya memerlukan penjumlahan.
        \item Dekripsi memerlukan satu perkalian modulo dan enam
              perbandingan \textit{greedy}.
    \end{itemize}
\end{frame}

% =====================================================================
\section{Steganografi: Menyembunyikan Keberadaan Pesan}
% =====================================================================

\begin{frame}{Gagasan Penyembunyian Informasi}
    \begin{block}{Apa yang disembunyikan}
        Kriptografi menyembunyikan \emph{isi} pesan. Steganografi
        menyembunyikan \kunci{keberadaan} pesan itu sendiri.
    \end{block}
    \vspace{0.1cm}
    \begin{exampleblock}{\textit{The Prisoner's Problem}}
        \footnotesize
        Alice dan Bob dipenjara dan hanya dapat berkomunikasi melalui Wendy.
        Mereka harus bertukar pesan tanpa membuat Wendy curiga bahwa ada pesan
        rahasia. Pesan yang terlihat biasa justru merupakan syarat
        keberhasilannya.
    \end{exampleblock}
    \begin{itemize}
        \item Pesan disisipkan ke dalam media pembawa yang wajar.
        \item Media pembawa dapat berupa teks, citra, audio, atau video.
        \item Peninjau tidak mengetahui bahwa ada pesan yang perlu dicari.
    \end{itemize}
\end{frame}

\begin{frame}{Terminologi Steganografi}
    \begin{itemize}
        \item \textbf{\textit{Embedded message}} --- pesan yang
              disembunyikan; dapat berupa teks, citra, audio, atau video.
        \item \textbf{\textit{Cover-object}} --- media yang dipakai untuk
              menyembunyikan pesan.
        \item \textbf{\textit{Stego-object}} --- media yang sudah berisi
              pesan tersembunyi.
        \item \textbf{\textit{Stego-key}} --- kunci yang dipakai untuk
              menyisipkan dan mengekstraksi pesan.
    \end{itemize}
    \vspace{0.2cm}
    \begin{columns}[T,onlytextwidth]
        \begin{column}{0.48\textwidth}
            \begin{block}{Sebelum}
                \centering
                \textit{cover-object} + pesan
            \end{block}
        \end{column}
        \begin{column}{0.48\textwidth}
            \begin{block}{Sesudah}
                \centering
                \textit{stego-object} yang tampak biasa
            \end{block}
        \end{column}
    \end{columns}
\end{frame}

\begin{frame}{Steganografi Bukan Pengganti Kriptografi}
    \begin{columns}[T,onlytextwidth]
        \begin{column}{0.54\textwidth}
            \footnotesize
            \begin{itemize}
                \item Keduanya saling melengkapi.
                \item Pesan dapat dienkripsi lebih dahulu, lalu cipherteksnya
                      disembunyikan di dalam media.
                \item Bila pesan itu ditemukan, penyerang masih harus
                      memecahkan ciphernya.
                \item Steganografi menurunkan peluang pesan terdeteksi;
                      kriptografi melindungi isinya bila terdeteksi.
            \end{itemize}
        \end{column}
        \begin{column}{0.42\textwidth}
            \begin{tikzpicture}[
                kotak/.style={draw, rounded corners, align=center,
                              font=\scriptsize, inner sep=3pt,
                              text width=2.6cm},
                panah/.style={-{Latex[length=1.8mm]}, thick, ubbgray},
            ]
                \node[kotak, fill=ubbblue!8] (p) at (0,0) {Pesan};
                \node[kotak, fill=ubbblue!8] (e) at (0,-1.1) {Enkripsi};
                \node[kotak, fill=ubbgreen!10] (s) at (0,-2.2)
                    {\textit{stego-object}};
                \draw[panah] (p) -- (e);
                \draw[panah] (e) -- (s);
                \node[font=\tiny, right=2pt of e] {cipherteks};
            \end{tikzpicture}
        \end{column}
    \end{columns}
\end{frame}

\begin{frame}{Kriteria Steganografi yang Baik}
    \begin{enumerate}
        \item \textbf{\textit{Imperceptible}} --- keberadaan pesan tidak dapat
              dipersepsi secara visual atau audial.
        \item \textbf{\textit{Fidelity}} --- kualitas media pembawa tidak jauh
              berubah.
        \item \textbf{\textit{Recovery}} --- pesan dapat diekstraksi kembali.
        \item \textbf{\textit{Capacity}} --- ukuran pesan yang disembunyikan
              sedapat mungkin besar.
    \end{enumerate}
    \vspace{0.15cm}
    \begin{alertblock}{Yang justru bukan isu utama}
        \footnotesize
        \textit{Robustness} atau ketahanan terhadap pengolahan media bukan
        prioritas utama dalam steganografi. Prioritasnya adalah pesan tidak
        terdeteksi dan tetap dapat dipulihkan.
    \end{alertblock}
\end{frame}

\begin{frame}{Media dan Ranah Operasi}
    \begin{columns}[T,onlytextwidth]
        \begin{column}{0.48\textwidth}
            \begin{block}{Media pembawa}
                \footnotesize
                \begin{itemize}
                    \item teks;
                    \item citra;
                    \item audio;
                    \item video.
                \end{itemize}
            \end{block}
        \end{column}
        \begin{column}{0.48\textwidth}
            \begin{block}{Ranah operasi}
                \footnotesize
                \begin{itemize}
                    \item \textbf{spasial/waktu}: mengubah langsung nilai byte
                          (intensitas piksel atau amplitudo);
                    \item \textbf{transformasi}: mengubah hasil transformasi
                          sinyal, misalnya ranah frekuensi.
                \end{itemize}
            \end{block}
        \end{column}
    \end{columns}
    \vspace{0.2cm}
    \begin{exampleblock}{Fokus pertemuan ini}
        \footnotesize
        Metode LSB termasuk ranah spasial. Ia mengubah bit kurang berarti pada
        nilai byte, baik pada citra maupun pada sampel audio.
    \end{exampleblock}
\end{frame}

\begin{frame}{Citra Digital}
    \begin{itemize}
        \item Citra terdiri atas sejumlah \textit{pixel}.
        \item Citra $1200 \times 1500$ memiliki
              $1200 \times 1500 = 1.800.000$ piksel.
        \item Setiap piksel panjangnya $n$ bit:
              \begin{itemize}
                  \item citra biner: $1$ bit per piksel;
                  \item citra \textit{grayscale}: $8$ bit per piksel;
                  \item citra \textit{true color}: $24$ bit per piksel.
              \end{itemize}
        \item Pada citra $24$ bit, satu piksel terdiri atas tiga komponen
              \textbf{RGB} (\textit{red}, \textit{green}, \textit{blue}).
        \item Setiap komponen menempati satu byte, sehingga satu piksel
              menyediakan tiga byte.
    \end{itemize}
    \vspace{0.1cm}
    \begin{exampleblock}{Konsekuensi}
        \footnotesize
        Setiap byte menyediakan satu bit LSB. Semakin banyak piksel, semakin
        besar pesan yang dapat disembunyikan.
    \end{exampleblock}
\end{frame}

\begin{frame}{Metode LSB}
    Di dalam setiap byte, bit tersusun dari yang paling berarti sampai yang
    paling kurang berarti:
    \[
        b_7b_6b_5b_4b_3b_2b_1b_0.
    \]
    \begin{itemize}
        \item \textbf{MSB} adalah $b_7$; \textbf{LSB} adalah $b_0$.
        \item Mengubah $b_0$ hanya mengubah nilai byte sebesar satu tingkat.
        \item Contoh: $1101000\underline{1}$ bernilai $209$; bila LSB diubah
              menjadi $0$, hasilnya $1101000\underline{0}$ bernilai $208$.
        \item Pergeseran satu dari $256$ tingkat warna tidak dapat dibedakan
              mata manusia.
    \end{itemize}
    \vspace{0.1cm}
    \begin{block}{Karena itu}
        \footnotesize
        LSB menjadi tempat alami untuk menyisipkan bit pesan tanpa mengubah
        persepsi atas citra secara keseluruhan.
    \end{block}
\end{frame}

\begin{frame}{Contoh LSB pada Satu Piksel}
    Tinjau satu piksel citra $24$ bit:
    \[
        [130,\;123,\;117]
        = \code{10000010}\;\code{01111011}\;\code{01110101}.
    \]
    Pesan yang disisipkan adalah tiga bit \code{101}.
    \begin{align*}
        \code{1000001}\underline{0} &\to \code{1000001}\underline{1}
            = 131, \\
        \code{0111011}\underline{1} &\to \code{0111011}\underline{0}
            = 122, \\
        \code{0111010}\underline{1} &\to \code{0111010}\underline{1}
            = 117.
    \end{align*}
    \begin{exampleblock}{Hasil}
        \footnotesize
        Piksel berubah dari $[130,123,117]$ menjadi $[131,122,117]$. Warna
        tetap tampak ``ungu''; perubahannya terlalu kecil untuk dilihat.
    \end{exampleblock}
\end{frame}

\begin{frame}{Penyisipan, Ekstraksi, dan Kapasitas}
    \begin{columns}[T,onlytextwidth]
        \begin{column}{0.50\textwidth}
            \begin{block}{Penyisipan dan ekstraksi}
                \footnotesize
                \begin{itemize}
                    \item Pesan dibaca bit demi bit.
                    \item Setiap bit menggantikan LSB satu byte.
                    \item Ekstraksi membaca byte, mengambil LSB-nya, lalu
                          merangkainya kembali menjadi bit pesan.
                \end{itemize}
            \end{block}
        \end{column}
        \begin{column}{0.46\textwidth}
            \begin{exampleblock}{Kapasitas}
                \footnotesize
                \begin{itemize}
                    \item \textit{Grayscale} $256\times256$:
                          $65536$ byte $\to$ $65536$ bit $=8$ KB.
                    \item Warna $24$ bit $256\times256$:
                          $196608$ byte $\to$ $196608$ bit $=24$ KB.
                \end{itemize}
            \end{exampleblock}
        \end{column}
    \end{columns}
    \vspace{0.15cm}
    \begin{alertblock}{Batasnya}
        \footnotesize
        Satu byte hanya menampung satu bit pesan. Kapasitas naik sebanding
        dengan jumlah byte media, bukan dengan ukuran berkas terkompresi.
    \end{alertblock}
\end{frame}

\begin{frame}{Varian Metode LSB}
    \begin{itemize}
        \item \textbf{Penyisipan berurutan} --- bit pesan menempati piksel
              pertama, kedua, dan seterusnya.
        \item \textbf{Penyisipan acak} --- posisi piksel dipilih oleh
              \textit{PRNG}; penerima membangkitkan urutan yang sama.
        \item \textbf{$m$-bit LSB} --- lebih dari satu bit LSB per byte
              dipakai, misalnya $b_1$ dan $b_0$.
        \item \textbf{Pra-enkripsi} --- bit pesan di-XOR dengan keystream
              sebelum disisipkan.
    \end{itemize}
    \vspace{0.15cm}
    \begin{alertblock}{\textit{Trade-off} yang harus diingat}
        \footnotesize
        Semakin banyak bit LSB yang dipakai, semakin besar kapasitas, tetapi
        semakin turun kualitas \textit{stego-image}. Bila posisi acak dipakai,
        diperlukan \textit{stego-key} untuk membangkitkan posisi tersebut.
    \end{alertblock}
\end{frame}

\begin{frame}{LSB pada Audio}
    \begin{columns}[T,onlytextwidth]
        \begin{column}{0.52\textwidth}
            \footnotesize
            \begin{itemize}
                \item Media yang dipakai adalah WAV tidak terkompresi.
                \item Berkas WAV tersusun atas \textit{RIFF header},
                      \textit{fmt chunk}, dan \textit{data chunk}.
                \item \textit{Data chunk} memuat sampel audio mentah.
                \item Pada PCM $16$ bit, satu sampel menempati $2$ byte.
                \item Stereo $16$ bit: satu frame berisi $4$ byte
                      (kiri dan kanan).
            \end{itemize}
        \end{column}
        \begin{column}{0.44\textwidth}
            \begin{block}{Penyisipan}
                \footnotesize
                \begin{itemize}
                    \item Satu bit pesan menggantikan LSB satu byte sampel.
                    \item Ekstraksi membaca LSB sampel dalam urutan yang sama.
                    \item Perubahan satu tingkat amplitudo tidak terdengar.
                \end{itemize}
            \end{block}
        \end{column}
    \end{columns}
    \vspace{0.15cm}
    \begin{alertblock}{Syarat praktis}
        \footnotesize
        Audio harus tidak terkompresi. Kompresi menghapus atau mengubah bit
        LSB, sehingga pesan tidak lagi dapat dipulihkan.
    \end{alertblock}
\end{frame}

% =====================================================================
\section{Steganalisis: Mendeteksi yang Tersembunyi}
% =====================================================================

\begin{frame}{Tujuan Steganalisis}
    \begin{block}{Tujuan utama}
        Menentukan apakah sebuah media yang dicurigai mengandung pesan
        tersembunyi atau tidak.
    \end{block}
    \vspace{0.15cm}
    \begin{itemize}
        \item Steganalisis dianggap berhasil bila keputusannya lebih baik
              daripada menerka secara acak.
        \item Tujuan minor: menentukan panjang pesan, tipe algoritma
              penyisipan, dan kunci yang dipakai.
    \end{itemize}
    \vspace{0.15cm}
    \begin{exampleblock}{Konteks forensik}
        \footnotesize
        Steganalisis dipakai di dalam \textit{forensic image analysis},
        khususnya \textit{content analysis}: mendeteksi pesan tersembunyi di
        dalam citra yang menjadi barang bukti.
    \end{exampleblock}
\end{frame}

\begin{frame}{\textit{Targeted} versus \textit{Blind} Steganalisis}
    \footnotesize
    \renewcommand{\arraystretch}{1.2}
    \begin{tabularx}{\linewidth}{@{}>{\raggedright\arraybackslash}p{2.6cm}
        >{\raggedright\arraybackslash}X
        >{\raggedright\arraybackslash}X@{}}
        \toprule
        Aspek & \textit{Targeted} & \textit{Blind} \\
        \midrule
        Cakupan & Algoritma steganografi tertentu, kadang terbatas pada format
                 media tertentu. & Sembarang algoritma dan sembarang format
                 media. \\
        \addlinespace
        Cara kerja & Mempelajari algoritma penyisipan, lalu mencari statistik
                    yang berubah setelah penyisipan. & Mempelajari perbedaan
                    statistik \textit{cover} dan \textit{stego}, biasanya
                    melalui pelatihan model \textit{machine learning}. \\
        \addlinespace
        Akurasi & Sangat akurat untuk sasaran yang tepat. & Kurang akurat
                  dibandingkan \textit{targeted}. \\
        \addlinespace
        Keluwesan & Tidak dapat diperluas ke algoritma atau format lain. &
                    Dapat diperluas ke algoritma lain. \\
        \bottomrule
    \end{tabularx}\par
    \vspace{0.3cm}
    \raggedright
    Keduanya menjawab pertanyaan yang sama dengan trade-off yang berbeda.
\end{frame}

\begin{frame}{Metode Steganalisis: Visual dan Statistik}
    \begin{columns}[T,onlytextwidth]
        \begin{column}{0.48\textwidth}
            \begin{block}{Serangan visual}
                \footnotesize
                \begin{itemize}
                    \item Khusus untuk \textit{stego-object} berupa citra.
                    \item Mengamati artefak secara kasat mata, lalu
                          membandingkannya dengan citra asli.
                    \item Bersifat subjektif.
                    \item Contoh: \textit{LSB plane attack} dan
                          \textit{filtered visual attack}
                          (\textit{enhanced LSB}).
                \end{itemize}
            \end{block}
        \end{column}
        \begin{column}{0.48\textwidth}
            \begin{block}{Serangan statistik}
                \footnotesize
                \begin{itemize}
                    \item Memakai analisis matematik atas citra.
                    \item Mencari perbedaan statistik antara \textit{cover}
                          dan \textit{stego}.
                    \item Penyisipan meninggalkan artefak yang dapat diukur.
                    \item Contoh: analisis histogram,
                          \textit{regular-singular} (RS),
                          \textit{chi-square}, dan
                          \textit{sample pair} (SP).
                \end{itemize}
            \end{block}
        \end{column}
    \end{columns}
    \vspace{0.1cm}
    \footnotesize
    Serangan visual banyak dipakai pada masa awal riset steganalisis;
    serangan statistik memberi ukuran yang lebih objektif.
\end{frame}

\begin{frame}{Batas Serangan Visual}
    \begin{exampleblock}{\textit{Enhanced LSB} pada citra kontras tinggi}
        \footnotesize
        Metode ini bekerja baik bila citra memiliki latar yang jelas atau
        perbedaan warna yang kontras antara latar dan objek utama. Artefak
        pesan menjadi mudah terlihat.
    \end{exampleblock}
    \vspace{0.1cm}
    \begin{alertblock}{Pada citra kontras rendah}
        \footnotesize
        Pada citra fotografi dengan kontras rendah, \textit{enhanced LSB}
        sering menyulitkan steganalisis. Pemeriksa sulit membedakan mana
        struktur gambar yang wajar dan mana artefak pesan rahasia.
    \end{alertblock}
    \vspace{0.1cm}
    \begin{block}{Pelajaran}
        \footnotesize
        Tidak ada satu metode yang cukup untuk semua citra. Pemeriksaan visual
        perlu dilengkapi metode statistik.
    \end{block}
\end{frame}

% =====================================================================
\section{Rangkuman dan Penutup}
% =====================================================================

\begin{frame}{Rangkuman}
    \footnotesize
    \begin{itemize}
        \item Knapsack mengubah bit plainteks menjadi pilihan objek dengan
              bobot total tertentu.
        \item Barisan \textit{superincreasing} dapat dipecahkan secara
              \textit{greedy} dalam waktu linear, sehingga tidak aman dipakai
              langsung.
        \item Merkle--Hellman menyembunyikannya melalui
              $w'_i=(w_i n)\bmod m$ dengan $m>\sum w_i$ dan
              $\mathrm{PBB}(n,m)=1$.
        \item Dekripsi memakai $C'=(C n^{-1})\bmod m$, lalu \textit{greedy}
              pada kunci privat.
        \item Contoh: $W_{priv}=\{2,3,6,13,27,52\}$, $m=105$, $n=31$,
              $n^{-1}=61$; blok \code{011000}, \code{110101}, \code{101110}
              menjadi $(174,280,333)$ dan pulih tepat.
        \item Shamir (1982) memecahkan Merkle--Hellman dasar dalam waktu
              polinomial dengan memanfaatkan struktur kuncinya.
        \item Steganografi menyembunyikan keberadaan pesan; LSB bekerja pada
              citra dan audio tidak terkompresi.
        \item Steganalisis dibedakan menjadi \textit{targeted} dan
              \textit{blind}, dengan metode visual dan statistik.
    \end{itemize}
\end{frame}

\begin{frame}{Yang Harus Diingat}
    \begin{columns}[T,onlytextwidth]
        \begin{column}{0.48\textwidth}
            \begin{block}{Knapsack}
                \footnotesize
                \begin{itemize}
                    \item Mudah dienkripsi, sukar didekripsi tanpa jebakan.
                    \item Jebakannya adalah barisan \textit{superincreasing}
                          di balik transformasi modulo.
                    \item Keamanannya runtuh karena strukturnya khusus, bukan
                          karena knapsack umum terpecahkan.
                \end{itemize}
            \end{block}
        \end{column}
        \begin{column}{0.48\textwidth}
            \begin{block}{Steganografi dan steganalisis}
                \footnotesize
                \begin{itemize}
                    \item LSB mengubah satu tingkat nilai byte; mata dan
                          telinga tidak menangkapnya.
                    \item Kapasitas terbatas oleh jumlah byte media.
                    \item Setiap penyisipan meninggalkan jejak yang dapat
                          dicari secara visual atau statistik.
                \end{itemize}
            \end{block}
        \end{column}
    \end{columns}
    \vspace{0.15cm}
    \begin{exampleblock}{Pertemuan berikutnya}
        \textbf{Minggu 15} --- Review CPMK-1 dan CPMK-2.
    \end{exampleblock}
\end{frame}

\begin{frame}{Penutup}
    \begin{block}{Tiga butir yang dibawa pulang}
        \footnotesize
        \begin{enumerate}
            \item Knapsack menunjukkan bahwa kunci publik dapat dibangun dari
                  persoalan yang sukar, asalkan pemilik kunci memegang jalan
                  pintas.
            \item Steganografi melengkapi kriptografi: isi pesan dilindungi,
                  keberadaannya disamarkan.
            \item Steganalisis mengingatkan bahwa penyembunyian selalu
                  meninggalkan jejak, sekecil apa pun perubahannya.
        \end{enumerate}
    \end{block}
    \vspace{0.2cm}
    \begin{exampleblock}{Pertemuan berikutnya}
        \textbf{Minggu 15} --- Review CPMK-1 dan CPMK-2: ringkasan seluruh
        algoritma dan konsep dari pertemuan 1 sampai 14.
    \end{exampleblock}
    \vspace{0.2cm}
    \centering
    \textit{Terima kasih.}
\end{frame}

\end{document}
